\section{Related work}
In this section, we will focus on summarizing relevant prior methods for deformable object manipulation. 

\textbf{Goal-conditioned manipulation} 
is particularly difficult for deformable objects due to their high degree-of-freedom and under-actuation \cite{7989247, 8575448,goal_transporter,ken_dense_descriptor,fabric_vsf_2020,abbeel_early}. Previous work studied tasks such as rope knot-tying \cite{ken_dense_descriptor} and fabric folding \cite{fabric_vsf_2020}. Early methods attempt to transfer demonstration trajectories to unseen configurations \cite{abbeel_early}. Recently, deep learning has been used extensively \cite{goal_transporter,ken_dense_descriptor,fabric_vsf_2020} to enable policy generalization to unseen objects or multiple tasks. However, these methods only consider quasi-static actions, where action effects are inferred mostly from object geometry alone. 
While quasi-static manipulations has shown to be sufficient to solve many tasks, the resulting systems are often slow and inefficient.
% In contrast, our method is able to reach goals far beyond the robot's reach range by exploiting the dynamics of novel ropes and clothes.

\textbf{Dynamic manipulation}
takes advantage of momentum (in addition to kinematic, static and quasi-static mechanisms) to increase load capacity or enable motions to have manipulands extend outside of the robot's nominal reach range \cite{dynamic_og,tossingbot,8344486,7139990,8280543,fast_hand,japan_shake}. Deformable objects pose a unique generalization challenge for dynamic manipulation policies due to complex physical properties and a particularly large sim2real gap.
%
Using reinforcement learning, Jangir et al. \cite{dyn_cloth_rl} proposes a system to dynamically fold and place cloth in various goal configurations. However, thousands of interactions are required for training. Similarly, Lim et al.\cite{casting} closes the sim2real gap by algorithmically tuning physics simulation parameters with real-world observations. But this process needs to be repeated for every unseen object instance. 
%
Leveraging visual feedback and domain randomization, Hietala et al \cite{finnish} is able to successfully transfer a manipulation policy trained in simulation to a real robotic setup. However, the policy is limited to a narrow range of goals such as folding clothes exactly in half. In comparison, our method is able to generalize to both unseen real-world object instances as well as wide variety of goals, making it a step closer to real-world applications.

\textbf{Trajectory optimization} uses analytical models \cite{fast_hand} or numerical simulation \cite{peter_model,bio_whip,optimal_planning,trajopt_berkeley,japan_rope} to generate mostly open-loop solutions for deformable object manipulation. Once a model has been designed, or automatically identified\cite{japan_rope}, the optimization problem can be solved using methods such as direct collocation \cite{trajopt_berkeley}, single shooting \cite{peter_model} or black-box optimization \cite{bio_whip}. These methods are generally capable of handling a wide range of goals and generalizing to many object instances, as long as a model is available for each new object. However, direct execution of optimal trajectory solutions on hardware (i.e., the OptSim method described later) does not work well for our tasks due to the large sim2real gap. In contrast, our method bridges this gap using feedback from previous trajectories without an explicit dynamics model.

\textbf{Iterative Learning Control (ILC)} leverages feedback from previous iterations to improve the accuracy of a repeatable system \cite{ilc_book, ilc_survey}. Most ILC algorithms tackle the task of trajectory tracking control, compensating for repeated disturbances against a reference trajectory. 
%However, when a forward dynamics model is not available, or when the the loss function is difficult to optimize over,  planning the optimal trajectory for a goal is a non-trivial problem on its own. 
However, when the objective function is non-convex and difficult to optimize (like the two tasks we investigate in this paper), the reference trajectory is hard to obtain. Hence, ILC can not be directly applied.
% However, since it is highly non-trivial to solve for the optimal trajectory for a given goal, ILC can not be directly applied. 
In addition, existing ILC algorithms often assumes known control direction or Jacobian which is difficult to obtain for our task. In contrast, our method is able to iteratively solve these tasks despite its complexity.
